\documentclass[a4paper,10pt]{article}
\usepackage{cw2026-report}

\newcommand{\ReportTitle}{Encoder odometry heading scale in 90\textdegree{} junction turns}
\newcommand{\ReportSubtitle}{Measured error, repeatability, and implications for the digital twin}

% EDIT TO ADD YOUR UNIQUE BRISTOL STUDENT ID NUMBER (7 digits).
\newcommand{\StudentID}{2071736}

\begin{document}
\ReportHeader

\begin{multicols*}{2}
\raggedcolumns

\section*{Subsystem and investigation aim}
Odometry is the only heading source available to our maze line-following
robot between junctions. The local model converts wheel counts to heading
using a wheel radius $r=16$~mm, wheel separation $b=90$~mm and
$N=358.3$~counts/rev (Pololu 3Pi+); the digital twin uses the same relationship. At every
90\textdegree{} junction the robot therefore relies on the encoder-to-heading
scale; an error in that scale accumulates along the route and makes the
twin's pose prediction diverge from the real robot.

This report asks how closely the encoder heading matches the physical heading
during a 90\textdegree{} in-place turn, whether the error is repeatable, and
what effective wheel separation explains it. If the nominal geometry were
correct, the encoder heading $\theta_{\mathrm{odom}}$ would be
90\textdegree{} with zero mean error; a systematic offset would instead
identify a geometry calibration error.

\section*{Measurement approach}
Firmware in \texttt{Controller.cpp} waits for 500~ms with no encoder change,
resets odometry, then drives both wheels at $\pm35$~PWM until the
differential count reaches 503.86~counts, the nominal 90\textdegree{} target
$\Delta N=\theta bN/(2\pi r)$. It then stops for 1.5~s. Ten trials (five left,
five right, alternating) were run on one A4 sheet taped to a desk, with the
robot centred on a printed cross and the same cable routing. Trial~5 was
rejected after the USB cable snagged; the 4~s safety timeout stopped it at
16.5\textdegree{}.

The independent reference is the onboard LSM6DS33 gyro (104~Hz,
$\pm245$~dps). Its $z$ rate is integrated and the bias is re-estimated during
each 500~ms quiet window. Per trial we compute $\theta_{\mathrm{odom}}$ at the
stop and after settling, $\theta_{\mathrm{true}}$ from the gyro, the
direction-corrected error
$e=\mathrm{sgn}(\theta)(\theta_{\mathrm{odom}}-\theta_{\mathrm{true}})$, and
$b_{\mathrm{eff}}=\Delta N\,(2\pi r/N)/|\theta_{\mathrm{true}}|$. The gyro
scale was checked against a printed protractor (agreement within about
3\textdegree{}); the remaining uncertainty is slip and coast.

% Replace all instructional text and captions with your own account.
% Keep one A4 page, with one or two plots. See README.md for guidance.
\ReportPlot{figures/fig_timeseries.pdf}{Heading during representative trials}
  {Encoder and gyro headings for one left (trial~1) and one right
   (trial~2) turn. The dashed line is the stop; afterwards the gyro coasts
   while the encoder changes only with the remaining counts.}
  {fig:timeseries}

\columnbreak

\section*{Results and interpretation}
Figure~\ref{fig:timeseries} shows representative left and right trials. The
two headings track during the turn but separate by the stop: in trial~1 the
encoder reports 90.6\textdegree{} while the gyro has reached
94.2\textdegree{}. The gyro then continues to change after the stop, whereas
the encoder only accumulates the remaining wheel counts.

\ReportPlot{figures/fig_summary.pdf}{Per-trial scale and final error}
  {(a) Effective separation from counts and gyro at the stop. (b) Final
   encoder under-report after settling. Circles: left; triangles: right;
   shaded band: rejected trial~5; solid/dashed: mean and nominal/zero.}
  {fig:summary}

Across the ten valid trials, $\theta_{\mathrm{odom}}$ at the stop was
$90.52\textdegree{}\pm0.25\textdegree{}$ and the gyro
$94.77\textdegree{}\pm1.1\textdegree{}$. The encoder under-reported the turn
in every trial: $4.25\textdegree{}\pm1.09\textdegree{}$ at the stop (95\% CI
3.47--5.03\textdegree{}) and $5.17\textdegree{}\pm1.75\textdegree{}$ after
settling (CI 3.92--6.42\textdegree{}). The implied separation was
$85.98\pm0.99$~mm (CI 85.27--86.69~mm), 4.47\% below the nominal 90~mm
(Figure~\ref{fig:summary}a). Rescaling the odometry by $90/85.98$ removes the
stop bias (mean $-0.01\textdegree{}$, SD $1.09\textdegree{}$); the final
residual is $-0.92\textdegree{}\pm1.76\textdegree{}$ because the coast
rotation is not counted. No significant left/right difference was found ($p=0.135$ at the stop,
$p=0.437$ after settling): the scale error is a direction-independent
geometry error with about 1\textdegree{} (1.2\%) random variation.

\section*{Implications for the digital twin}
The twin instantiates the same model with $r=16$~mm and $b=90$~mm
(\texttt{DigitalTwin/Simulator.pde}), using $\Delta\theta=(d_R-d_L)/b$. Setting $b=86.0$~mm, or scaling $\theta$ by
1.047, removes the 4.5\% bias at the stop. With $b=90$~mm, four successive
90\textdegree{} turns accumulate about 17\textdegree{} of heading error at the
stop (about 21\textdegree{} after coast); the simulated robot would predict a
different corridor from the real one. The twin should also represent the short post-stop coast (mean 1.3\textdegree{}, up to
3.3\textdegree{}); integrating counts alone under-predicts the final heading
by about 1\textdegree{} per turn. A check is five fresh 90\textdegree{} turns
with $b=86.0$~mm; a mean residual above 1\textdegree{} or an SD above
2\textdegree{} would require revisiting the calibration or coast model. The
result applies to in-place turns at PWM~35 on paper; other conditions need a
separate check.

\EvidenceNote{Code: \texttt{MyFirstRobot/Controller.*}, \texttt{MyFirstRobot.ino},
\texttt{OdometryModel.*}. Data:
\path{test_materials/turn_capture_main.txt} (trial~5 rejected);
\path{results/turn_trial_summary.csv}. Analysis:
\path{analysis/analyse_turn_trials.py}.}

\end{multicols*}
\end{document}
